% ===================== ALCANCE =====================
\section{Alcance y Limitaciones}
\label{sec:alcance}

\subsection{Alcance del Estudio}

\begin{table}[htbp]
    \centering
    \caption{Elementos incluidos en el alcance del an\'{a}lisis.}
    \label{tab:alcance_incluido}
    \setcounter{itemcount}{0}
    \begin{tabularx}{\textwidth}{@{}NX@{}}
        \toprule
        \multicolumn{1}{c}{\textbf{Item}} & \textbf{Descripcion} \\
        \midrule
        & Modelo de elementos finitos de la l\'{i}nea SIM-008 (job CAESAR P2603-PR-SIM-008), desde la salida de la lavadora DDW hasta la boquilla de succi\'{o}n de la bomba de recirculaci\'{o}n 1, generado desde el isom\'{e}trico PCF106. \\
        & Evaluaci\'{o}n de esfuerzos seg\'{u}n ASME B31.3-2016 para nueve casos de carga: dos de operaci\'{o}n (OPE), dos sostenidos alternos (Alt-SUS), dos sostenidos (SUS) y tres de expansi\'{o}n (EXP). \\
        & Reporte de desplazamientos nodales, esfuerzos por elemento y cargas en soportes para los casos de carga relevantes. \\
        & Verificaci\'{o}n de cumplimiento de c\'{o}digo (\textit{code compliance}) y elaboraci\'{o}n del presente informe t\'{e}cnico con c\'{o}digo de documento I24.104-PP30-PP01-P-DOCS-183. \\
        \bottomrule
    \end{tabularx}
\end{table}

\subsection{Limitaciones del Estudio}

\begin{table}[htbp]
    \centering
    \caption{Elementos excluidos del alcance del an\'{a}lisis.}
    \label{tab:alcance_excluido}
    \setcounter{itemcount}{0}
    \begin{tabularx}{\textwidth}{@{}NX@{}}
        \toprule
        \multicolumn{1}{c}{\textbf{Item}} & \textbf{Descripcion} \\
        \midrule
        & Cargas ocasionales de viento y sismo: no reportadas en el modelo; su evaluaci\'{o}n, de ser requerida, corresponde a un an\'{a}lisis posterior. \\
        & Cargas en la boquilla de succi\'{o}n de la bomba frente a criterios de fabricante (p.\,ej. API 610): el modelo reporta cargas en el extremo de succi\'{o}n, pero su comparaci\'{o}n con los admisibles del equipo requiere los datos del fabricante, no disponibles a la fecha. \\
        & Fatiga detallada por ciclos de arranque/parada: el criterio de expansi\'{o}n de B31.3 con factor $f$ cubre el rango t\'{e}rmico; no se ejecut\'{o} an\'{a}lisis de fatiga espec\'{i}fico. \\
        & Verificaci\'{o}n en campo de la soporter\'{i}a: las propiedades de los soportes corresponden a las definidas en el isom\'{e}trico de dise\~{n}o. \\
        \bottomrule
    \end{tabularx}
\end{table}
